\subsection{6.3 Metodologia de Avaliação Causal}

Para avaliar formalmente o impacto causal do transtorno mental no desfecho de lesão autoprovocada, implementou-se um Modelo Estrutural Causal (SCM) baseado no grafo direcionado acíclico (DAG) de especialistas estabelecido anteriormente. Utilizando o framework \textit{DoWhy}, o primeiro passo consistiu na identificação do efeito causal para a Hipótese 5.1, definindo a presença de transtorno mental (\texttt{TRAN\_MENT}) como a variável de tratamento e a ocorrência de lesão autoprovocada (\texttt{LES\_AUTOP}) como a variável alvo (desfecho).

A análise topológica do grafo determinou que o efeito é identificável por meio do Critério \textit{Backdoor} (Caminho Secundário). Para isolar o efeito do tratamento e neutralizar a ação de variáveis confundidoras, o algoritmo determinou a necessidade de condicionamento e ajuste estatístico sobre um conjunto específico de covariáveis observadas no dataset do SINAN: a faixa etária (\texttt{NU\_IDADE\_N}), a situação conjugal (\texttt{SIT\_CONJUG}) e o histórico de recorrência de violência (\texttt{OUT\_VEZES}). Sob a premissa de ausência de confundimento não observado (\textit{Unconfoundedness Assumption}), o bloqueio dessas vias garante que a associação remanescente entre o tratamento e o desfecho reflita puramente uma relação de causa e efeito.

\subsection{6.4 Resultados da Estimação}

Após a execução da limpeza rigorosa e do isolamento estrito das matrizes matemáticas do modelo, a base de dados nacional de 2024 totalizou uma amostra qualificada de 96.193 observações biunívocas. Dado que a variável de desfecho (\texttt{LES\_AUTOP}) é estritamente binária (representando a ocorrência ou não de lesão autoprovocada), utilizou-se um Modelo Linear Generalizado (GLM) com família Binomial, equivalente a uma abordagem de Regressão Logística, para computar o Efeito Causal Médio (ATE).

Os resultados quantitativos da estimação são apresentados a seguir:
\begin{itemize}
    \item \textbf{Efeito Causal Médio (ATE) Estimado:} 0,156084
    \item \textbf{Intervalo de Confiança de 95\%:} (0,149822, 0,162096)
\end{itemize}

O valor obtido para o ATE indica que, mantendo todas as variáveis de confusão controladas e constantes, a presença diagnosticada de um transtorno mental (\texttt{TRAN\_MENT}) introduz um incremento líquido de \textbf{15,61\%} na probabilidade absoluta de o indivíduo apresentar um comportamento de lesão autoprovocada. O intervalo de confiança extremamente estreito evidencia a alta precisão estatística da estimativa, decorrente do robusto volume amostral processado.

\subsection{6.5 Testes de Robustez (Refutação)}

Para garantir que o efeito causal estimado não seja decorrente de artefatos estatísticos, correlações espúrias ou falhas de especificação do DAG, o modelo foi submetido a dois testes rigorosos de refutação baseados em simulações contrafactuais.

\begin{enumerate}
    \item \textbf{Adição de Causa Comum Aleatória (\textit{Random Common Cause}):} Este teste insere uma variável independente artificial (ruído gaussiano) no conjunto de dados para atuar como um confundidor fictício. Caso o modelo original dependesse de uma especificação frágil, a inclusão desse ruído deslocaria significativamente o efeito. A reestimação do modelo sob essa condição resultou em um novo efeito de \textbf{0,156083}, demonstrando estabilidade quase absoluta. O valor-p obtido foi de \textbf{0,76}, confirmando que a introdução do novo termo não possui impacto estatisticamente significante sobre a estimativa original.
    \item \textbf{Tratamento Placebo (\textit{Placebo Treatment}):} Este teste substitui a variável real de tratamento (\texttt{TRAN\_MENT}) por uma variável puramente aleatória, obtida por meio da permutação dos dados originais. Como o ``falso tratamento'' não possui conexão causal real com o desfecho, o efeito estimado deve colapsar. A execução do teste confirmou essa premissa: o novo efeito caiu para \textbf{$8,22 \times 10^{-14}$} (zero do ponto de vista prático), com um valor-p de \textbf{0,0}. A rejeição da hipótese nula neste cenário valida o modelo, assegurando que o efeito de 15,61\% observado anteriormente é de fato gerado pela variação de \texttt{TRAN\_MENT}.
\end{enumerate}

\section{7. Conclusão}

Este trabalho apresentou uma abordagem estruturada de Aprendizado de Máquina Causal para investigar os determinantes e fatores de risco associados à violência autoprovocada no Brasil, utilizando dados epidemiológicos do SINAN/Datasus do ano de 2024. Diante das limitações biológicas e lógicas evidenciadas pelos algoritmos de descoberta causal puramente orientados a dados — os quais geraram inconsistências severas, como a inversão temporal entre o desfecho e construtos de identidade de gênero —, a incorporação de restrições baseadas em conhecimento de domínio mostrou-se um passo metodológico indispensável.

A modelagem causal e os experimentos subsequentes confirmaram com robustez estatística a Hipótese 5.1 formulada neste projeto. O diagnóstico ou presença de transtornos mentais atua como um forte determinante causal em direção à lesão autoprovocada, elevando o risco absoluto do desfecho em 15,61\%. Os testes contrafactuais de refutação blindaram as conclusões, rejeitando a existência de vieses por confundimento não observado ou aleatoriedade nos dados.

Do ponto de vista de formulação de políticas públicas e intervenções em saúde mental, esses achados fornecem um forte argumento quantitativo. Eles demonstram que estratégias direcionadas ao suporte clínico, acompanhamento médico e mitigação de sofrimentos psíquicos crônicos e transtornos comportamentais possuem um canal de causalidade direto e mensurável, capaz de reduzir significativamente as taxas de violência autoprovocada e tentativas de autoextermínio na população.
